% GENERATED FILE. Edit canonical lessons, then run npm run generate:books.
% canonical-source-hash: fnv1a64:9b26ff9f3a7d2537
% canonical-lessons: ZH-C182-fanguan, ZH-C182-lengcai, ZH-W182-leng-cold, ZH-C182-xiaochi, ZH-C182-jiachangcai, ZH-C182-hefan

\chapter{Eating Out}
\label{ch:zh-eating-out}
\hlchaptermodality{182}

\begin{chapteropening}
\textbf{By the end of this chapter:} \emph{I can talk about restaurants, cold starters, snacks, home cooking and boxed meals.}

Complete the last lesson of chapter 182: I can talk about restaurants, cold starters, snacks, home cooking and boxed meals.
\end{chapteropening}

\section[fànguǎn]{fànguǎn — restaurant}

\label{lesson:ZH-C182-fanguan}



\begin{quote}
Before the new one: say the Mandarin for a bathroom, then the Mandarin for a curtain.
\end{quote}

\subsection*{You'll want to know: fànguǎn}
\emph{fànguǎn} — \textquotedblleft{}restaurant\textquotedblright{}.

\begin{quote}
\emph{wǒ jiā hòubian yǒu yí ge fànguǎn} — There's a restaurant behind our place.
\end{quote}

Its characters come later; for now it is a word for the ear.

\begin{cousinweb}
\emph{fànguǎn} is built from \zh{饭} 'meal' and \emph{guǎn} 'lodging house', both inherited from Middle Chinese, \emph{bjonH} and \emph{kwanH} (Guangyun, in Baxter's transcription), though the falling tone of \emph{kwanH} would regularly give guàn.
\end{cousinweb}

\begin{grammarlens}[title={What you've built}]
One more word for asking for what I want.
\end{grammarlens}

\subsection*{Your turn}
\begin{itemize}
\raggedright
  \item \emph{Say it:} \emph{fànguǎn}
  \item \emph{Say it:} \emph{fànguǎn}, once more
  \item \emph{Say it:} \emph{wǒ jiā hòubian yǒu yí ge fànguǎn} — There's a restaurant behind our place.
  \item \emph{Recall:} say the Mandarin for a bathroom, then the Mandarin for a curtain, then say \emph{fànguǎn} again
\end{itemize}

\begin{tcolorbox}[breakable,colback=teal!4,colframe=teal!35!black,title={Before you move on}]
What does fànguǎn mean? (\textquotedblleft{}Restaurant\textquotedblright{}.) Say it once more.
\end{tcolorbox}
\section[\texorpdfstring{\zh{冷菜} (lěngcài)}{冷菜 (lěngcài)}]{\zh{冷菜} (lěngcài) — cold starter; cold dishes (on a menu)}

\label{lesson:ZH-C182-lengcai}



\begin{quote}
Before the new one: say the Mandarin for a curtain, then the Mandarin for a restaurant.
\end{quote}

\subsection*{You'll want to know: \zh{冷菜}}
\textbf{\zh{冷菜}} — \emph{lěngcài} — \textquotedblleft{}cold starter; cold dishes (on a menu)\textquotedblright{}.

\begin{quote}
\textbf{\zh{我们要两个冷菜}} \emph{wǒmen yào liǎng ge lěngcài} — We'd like two cold starters.
\end{quote}

The new character, in pieces: \zh{冷} — seven strokes: the ice piece on the left (a dot, then a stroke rising to the right), and on the right \zh{今} as you write it, with one more dot under it. The next lesson puts it in your hand.

\begin{cousinweb}
\textbf{\zh{冷菜}} (\emph{lěngcài}) is built from \zh{冷} 'cold' and \zh{菜} 'dish', both inherited from Middle Chinese, \emph{lengX} and \emph{tshojH} (Guangyun, in Baxter's transcription), which glosses the first 'cold' and calls the second any plant fit to eat; 'a dish' is the later sense. The dishes a meal starts with, served cold.
\end{cousinweb}

\begin{grammarlens}[title={What you've built}]
One more word for asking for what I want.
\end{grammarlens}

\subsection*{Your turn}
\begin{itemize}
\raggedright
  \item \emph{Say it:} \emph{lěngcài}
  \item \emph{Say it:} \emph{lěngcài}, once more
  \item \emph{Say it:} \emph{wǒmen yào liǎng ge lěngcài} — We'd like two cold starters.
  \item \emph{Recall:} say the Mandarin for a curtain, then the Mandarin for a restaurant, then say \emph{lěngcài} again
\end{itemize}

\begin{tcolorbox}[breakable,colback=teal!4,colframe=teal!35!black,title={Before you move on}]
What does \zh{冷菜} mean? (\textquotedblleft{}Cold starter; cold dishes (on a menu)\textquotedblright{}.) Say it once more.
\end{tcolorbox}
\section[\texorpdfstring{\zh{冷} (lěng)}{冷 (lěng)}]{\zh{冷} — the ice piece, then \zh{今} and a dot}

\label{lesson:ZH-W182-leng-cold}



\begin{quote}
The word first, then the character inside it.

\begin{itemize}
\raggedright
  \item \emph{Recall:} say \emph{lěngcài} — \textbf{R1}, one lesson back, before it decays
\end{itemize}

\begin{quote}
\zh{今} \zh{菜}
\end{quote}

\emph{Read it:} each one aloud, with nothing said first — \emph{jīn}, \emph{cài}
\end{quote}

\subsection*{Script — the ice piece, then \zh{今} and a dot}
\hlblockfigure{\detokenize{figures/ZH-W182-leng-cold-filmstrip.pdf}}{How \zh{冷} is written, stroke by stroke}

\begin{quote}
\zh{冷}
\end{quote}

\textbf{Two pieces, left before right.} Count the pieces, not the strokes.

\textbf{The ice piece, two strokes.} A dot. Then, under it, a short stroke rising to the upper right.

\textbf{The right piece, five strokes: \zh{今}, then a dot.} \zh{今} as you write it: a stroke falling to the left, one falling to the right, a dot under the roof, and a horizontal that turns and falls to the left. Last, one more dot under it.

2 + 4 + 1 = seven.

\subsection*{Your turn}
\begin{itemize}
\raggedright
  \item \emph{Write it:} \zh{冷}, all seven, with the model in front of you
  \item \emph{Cover:} write it again from memory, and count the pieces, not the strokes
\end{itemize}

\begin{tcolorbox}[breakable,colback=teal!4,colframe=teal!35!black,title={Before you move on}]
How many strokes in \zh{冷}? (\textbf{Seven}.) How many in the ice piece? (\textbf{Two: a dot, then a rising stroke}.) Which character makes most of the right-hand piece? (\textbf{\zh{今}, with one more dot}.)

Source: \href{https://raw.githubusercontent.com/chanind/hanzi-writer-data/68d10a4b21150cae5e1ebbd223eed289cf32d90c/data/\%E5\%86\%B7.json}{Hanzi Writer Data \zh{冷}, pinned stroke paths and medians}.
\end{tcolorbox}
\section[\texorpdfstring{\zh{小吃} (xiǎochī)}{小吃 (xiǎochī)}]{\zh{小吃} (xiǎochī) — snack; street food}

\label{lesson:ZH-C182-xiaochi}



\begin{quote}
Before the new one: say the Mandarin for a restaurant, then the Mandarin for a cold starter.
\end{quote}

\subsection*{You'll want to know: \zh{小吃}}
\textbf{\zh{小吃}} — \emph{xiǎochī} — \textquotedblleft{}snack; street food\textquotedblright{}.

\begin{quote}
\textbf{\zh{中午没有时间我们吃小吃}} \emph{zhōngwǔ méiyǒu shíjiān wǒmen chī xiǎochī} — There's no time at lunch, so we'll grab a snack.
\end{quote}

\begin{cousinweb}
\textbf{\zh{小吃}} (\emph{xiǎochī}) is built from \zh{小} 'small' and \zh{吃} 'to eat': a small eat. The first is inherited from Middle Chinese \emph{sjewX} (Guangyun, in Baxter's transcription); for the second the dictionary has \emph{khek}, glossed 'to eat', under another way of writing it, though \emph{khek} would regularly give qì.
\end{cousinweb}

\begin{grammarlens}[title={What you've built}]
One more word for asking for what I want.
\end{grammarlens}

\subsection*{Your turn}
\begin{itemize}
\raggedright
  \item \emph{Say it:} \emph{xiǎochī}
  \item \emph{Say it:} \emph{xiǎochī}, once more
  \item \emph{Say it:} \emph{zhōngwǔ méiyǒu shíjiān wǒmen chī xiǎochī} — There's no time at lunch, so we'll grab a snack.
  \item \emph{Recall:} say the Mandarin for a restaurant, then the Mandarin for a cold starter, then say \emph{xiǎochī} again
\end{itemize}

\begin{tcolorbox}[breakable,colback=teal!4,colframe=teal!35!black,title={Before you move on}]
What does \zh{小吃} mean? (\textquotedblleft{}Snack; street food\textquotedblright{}.) Say it once more.
\end{tcolorbox}
\section[\texorpdfstring{\zh{家常菜} (jiāchángcài)}{家常菜 (jiāchángcài)}]{\zh{家常菜} (jiāchángcài) — home cooking; everyday dishes}

\label{lesson:ZH-C182-jiachangcai}



\begin{quote}
Before the new one: say the Mandarin for a cold starter, then the Mandarin for a snack.
\end{quote}

\subsection*{You'll want to know: \zh{家常菜}}
\textbf{\zh{家常菜}} — \emph{jiāchángcài} — \textquotedblleft{}home cooking; everyday dishes\textquotedblright{}.

\begin{quote}
\textbf{\zh{我们去对面吃家常菜}} \emph{wǒmen qù duìmiàn chī jiāchángcài} — Let's go across the road for some home cooking.
\end{quote}

\begin{cousinweb}
\textbf{\zh{家常菜}} (\emph{jiāchángcài}) is built from \zh{家} 'home', \zh{常} 'everyday' and \zh{菜} 'dish', all three inherited from Middle Chinese, \emph{kae}, \emph{dzyang} and \emph{tshojH} (Guangyun, in Baxter's transcription): the dishes a home makes every day. The everyday sense of the second is not given there.
\end{cousinweb}

\begin{grammarlens}[title={What you've built}]
One more word for asking for what I want.
\end{grammarlens}

\subsection*{Your turn}
\begin{itemize}
\raggedright
  \item \emph{Say it:} \emph{jiāchángcài}
  \item \emph{Say it:} \emph{jiāchángcài}, once more
  \item \emph{Say it:} \emph{wǒmen qù duìmiàn chī jiāchángcài} — Let's go across the road for some home cooking.
  \item \emph{Recall:} say the Mandarin for a cold starter, then the Mandarin for a snack, then say \emph{jiāchángcài} again
\end{itemize}

\begin{tcolorbox}[breakable,colback=teal!4,colframe=teal!35!black,title={Before you move on}]
What does \zh{家常菜} mean? (\textquotedblleft{}Home cooking; everyday dishes\textquotedblright{}.) Say it once more.
\end{tcolorbox}
\section[héfàn]{héfàn — boxed meal; packed lunch}

\label{lesson:ZH-C182-hefan}



\begin{quote}
Before the new one: say the Mandarin for a snack, then the Mandarin for home cooking.
\end{quote}

\subsection*{You'll want to know: héfàn}
\emph{héfàn} — \textquotedblleft{}boxed meal; packed lunch\textquotedblright{}.

\begin{quote}
\emph{wǒmen zhōngwǔ chī héfàn} — We have boxed lunches at midday.
\end{quote}

Its characters come later; for now it is a word for the ear.

\begin{cousinweb}
\emph{héfàn} is built from \emph{hé} 'box' and \zh{饭} 'meal', both inherited from Middle Chinese, \emph{hop} and \emph{bjonH} (Guangyun, in Baxter's transcription), which glosses the first as a dish with a lid.
\end{cousinweb}

\begin{grammarlens}[title={What you've built}]
One more word for asking for what I want.
\end{grammarlens}

\subsection*{Your turn}
\begin{itemize}
\raggedright
  \item \emph{Say it:} \emph{héfàn}
  \item \emph{Say it:} \emph{héfàn}, once more
  \item \emph{Say it:} \emph{wǒmen zhōngwǔ chī héfàn} — We have boxed lunches at midday.
  \item \emph{Recall:} say the Mandarin for a snack, then the Mandarin for home cooking, then say \emph{héfàn} again
\end{itemize}

\begin{tcolorbox}[breakable,colback=teal!4,colframe=teal!35!black,title={Before you move on}]
What does héfàn mean? (\textquotedblleft{}Boxed meal; packed lunch\textquotedblright{}.) Say it once more.
\end{tcolorbox}
